% APPENDIX 1. The systematic member-by-member specification. The Theory section carries the
% narrative (two closures, the ladder, the two structural differences) and points here for the
% formulas; this file carries the complete cycle with the three switches, so that a reader can
% reconstruct any member without reading the code.
% src: theory/macroir/docs/Macro_IR/macroir_derivation.tex (the canonical derivation: bilinear
%      tilde 740-768 and its closed form 880-941; vector tilde 1148-1165)
% src: legacy/qmodel.h:4498-4602 (the implemented cycle: the av=0 pre-propagation 4498-4512,
%      the two gSg forms 4525-4537, the variance-form switch 4543-4576, the two gains 4594-4602)
% NOT sourced from theory/macroir/docs/Macro_IR/MacroIR_tilde.md, which its own header marks as
% predating the corrected normalization.

% \appendix and the A-numbering are declared once, in 08_appendix_derivation.tex, which is input
% before this file. Redeclaring them here would restart the counter and give two equations called A1.
\begin{appendixbox}
\label{app:members}
\section*{The interval update, member by member}

Appendix~\ref{app:derivation} gives the cycle for the boundary-conditioned member and names the two
places where the others differ. This appendix writes every member out in one template, so
that each is fixed by three switches on a single set of equations.

\subsection*{Objects computed once per interval}

All members share the quantities built from the rate matrix $\mathbf{Q}$ and the recording
filter for a window of length $\Delta$: the transition matrix $\mathbf{P}(\Delta)$; the
boundary-conditioned single-channel moments $\bar{\gamma}_{i\to j}$ and
$\bar{v}_{i\to j}$; the start-conditioned mean conductance $\bar{\bm{\gamma}}_0$ of
Eq.~\ref{eq:gammabar}; the start-conditioned interval variance
$\mathbf{w}$; and, for the members that ignore the acquisition
averaging, the instantaneous per-state conductance vector $\bm{\gamma}$. Two matrices
collect the boundary weights,
\begin{equation}
  G_{i\to j} \;=\; \bar{\gamma}_{i\to j}\,P_{i\to j}(\Delta),
  \qquad
  H_{i\to j} \;=\; \bar{\gamma}^2_{i\to j}\,P_{i\to j}(\Delta).
  \label{eq:GH}
\end{equation}

\subsection*{The tilde operator}

Over a bold symbol, the tilde marks a contraction taken on the $K^2$ boundary pairs $(i\to j)$
rather than on the $K$ states. (Over a scalar it means something else, a quantity in the natural
units of the model, as in $\widetilde{\Delta}$ and $\widetilde{S}$; the two uses are disjoint and
never meet in one expression.) It is not a new operator: its two cases are the two contractions the update of
Eq.~\ref{eq:macror} already has, $\bm{\gamma}^\top\bm{\Sigma}\bm{\gamma}$ and
$\bm{\Sigma}\bm{\gamma}$, evaluated there with the boundary objects substituted. Both are written out in Appendix~\ref{app:derivation}, the
bilinear one as the scalar of Eq.~\ref{eq:tilde} and the vector one as the gain
$\mathbf{g}_{0\to\Delta}=\widetilde{\bm{\gamma}^\top\bm{\Sigma}}$ of Eq.~\ref{eq:vector-tilde}, and both
collapse to $K\times K$ arithmetic once the boundary covariance Eq.~\ref{eq:sig-bnd} is
substituted, so the $K^2\times K^2$ object is never built. The matrices $\mathbf{G}$ and
$\mathbf{H}$ of Eq.~\ref{eq:GH} are what tells the two apart, one power of $\bar{\gamma}$
against two.
% MOVED 2026-08-06. The vector tilde was stated here and only referred to in words in the body,
% which sent a reader looking for the gain out of the section that introduces it. The equation now
% lives at 02_theory.tex next to Eq. tilde and keeps the label eq:vector-tilde, so every reference
% here still resolves; this paragraph points at it rather than restating it, which is already how
% it treats the bilinear case.
% AMENDED the same day, with the Theory section's rewrite onto Luciano's own derivation. The
% paragraph opened "The tilde is a lift, a modulation and a collapse", which promised a piece of
% apparatus where there is only MacroR's own quadratic form and gain read on the pairs. Saying that
% plainly costs nothing and removes the reader's suspicion that a K^2 machine is being hidden.

\subsection*{One cycle, three switches}

Every member on the lattice runs the same four steps on each window: propagate the occupancy
Gaussian (Eqs.~\ref{eq:mu-prop} and~\ref{eq:sig-prop}), predict the observation
(Eqs.~\ref{eq:ypred} and~\ref{eq:pred-var}), condition on the recorded average
(Eqs.~\ref{eq:mu-post} and~\ref{eq:sig-post}), and either carry the posterior to the next
window or discard it. Three switches fix which member is being run.

\paragraph{Switch 1: the window.} It has three settings, an instant, the state at the
window's start, and the states at both of its ends, and it selects the conductance object,
and with it the second term of the predictive variance and the gain. At the instantaneous
setting the observation is treated as a sample taken at the centre of
the window: the occupancy is propagated there first, by $\mathbf{P}(\Delta/2)$, the conductance
is $\bm{\gamma}$, and, with $\bm{\Sigma}$ read in that mid-window frame,
\begin{equation}
  \widetilde{\bm{\gamma}^\top\bm{\Sigma}\bm{\gamma}} \;\to\; \bm{\gamma}^\top\bm{\Sigma}\bm{\gamma},
  \qquad
  \mathbf{g} \;\to\; \mathbf{P}(\Delta/2)^\top\bm{\Sigma}\,\bm{\gamma}.
  \label{eq:av0}
\end{equation}
At the start-conditioned setting the same two expressions hold with $\bar{\bm{\gamma}}_0$ in
place of $\bm{\gamma}$ and the full $\mathbf{P}(\Delta)$ in the gain, since the conductance is
conditioned on the state at the window's start and the whole window remains to be crossed;
the contraction is still an ordinary quadratic form on the $K$ states.
The trailing propagator in the gain is not decoration. It carries the update from the frame in
which the observation was formed to the frame the next window starts in, which is the end of
the current window, so it is the remaining half at the instantaneous setting and the whole
window at the start-conditioned one. Omitting it, or using the full window where half is called for, leaves the
rank-one down-date in the wrong frame, and the symptom is specific and was observed: a spike in
the distortion-induced bias of the recursive instantaneous member at long intervals and large
channel counts.
% CORRECTED 2026-08-04. Eq. av0 read P(Delta) in the gain and named no frame for Sigma. Both are
% wrong for av=0: the implementation propagates to the window MIDPOINT and the trailing factor is
% P_half. Source, with the table that makes the three cases explicit and the historical failure it
% caused: docs/math/trust_coefficient.md:24-40 ("| 0 | X_mid via instantaneous g | g^T . Sigma_mid .
% P_half |", and "Without it the down-date is in the wrong frame - manifested historically as a
% Distortion-Induced-Bias spike in macro_R at long intervals / large Num_ch (fixed 2026-05-10)").
% Found by a replication audit, which noticed the appendix and that file disagree.
% CHECK BEFORE SUBMISSION: whether Delta/2 is exactly what the code uses, or whether P_half is
% defined against a different reference point (legacy/qmodel.h:4508-4515).
% FIXED 2026-08-06. The line above used to end "(legacy/qmodel.h:4508-4515). At $\mathrm{av}=2$ the",
% so the opening of the next body sentence was inside the comment and the printed paragraph ran
% "... at long intervals and large channel counts. conductance is conditioned on both endpoints".
% Same failure as the one repaired at 02_theory.tex:82 on the same day, in a different file: a
% sentence restarted on a commented line. Whenever a comment block is inserted mid-paragraph here,
% check that its last line carries no body text.
At the boundary setting the
conductance is conditioned on both endpoints and the two expressions become the tildes,
Eq.~\ref{eq:tilde} and Eq.~\ref{eq:vector-tilde}. The step from the start state to the
boundary therefore does two things at once: it replaces $\bm{\Sigma}_0$ by
$\bm{\Sigma}_0-\mathrm{diag}(\bm{\mu}_0)$ and adds the boundary term ($\bm{\mu}_0^\top
\mathbf{H}\mathbf{1}$ in the variance, $\mathbf{G}^\top\bm{\mu}_0$ in the gain). The
subtracted diagonal is not a second approximation: the within-channel part it removes is
already inside the boundary term, and dropping the subtraction would count it twice.

\paragraph{Switch 2: recursion.} A recursive member carries
$(\bm{\mu}^{\mathrm{post}}, \bm{\Sigma}^{\mathrm{post}})$ into the next window. A
non-recursive one discards it and propagates the next window's predictive from the initial
condition, so its log-likelihood factorizes over windows. The switch changes no formula
above; it changes only what is fed to the next one.

\paragraph{Switch 3: the interval variance, total or residual.} It selects the third term of
Eq.~\ref{eq:pred-var}, which is present only when the acquisition averaging is modelled at
all, that is at either of the two interval settings of switch~1:
\begin{equation}
  \bigl(\mathbf{w}\bigr)_{i} =
  \begin{cases}
    \sum_{j} P_{i\to j}(\Delta)\bigl[\bar{v}_{i\to j}
      + \bigl(\bar{\gamma}_{i\to j}-\bar{\gamma}_{i\to\cdot}\bigr)^2\bigr]
      & \text{total,}\\[4pt]
    \sum_{j} P_{i\to j}(\Delta)\,\bar{v}_{i\to j}
      & \text{residual.}
  \end{cases}
  \label{eq:variance-form}
\end{equation}
The two differ by the spread of the interval mean across end states, which is
Eq.~\ref{eq:mr-vr}. The switch is free only where the window is conditioned on the start
state alone. With both endpoints conditioned on, the
residual form is forced whatever the flag says, for the same reason the diagonal is
subtracted in switch~1: the between-endpoint spread is already carried by the boundary term,
so the total form would count it twice. % src: legacy/qmodel.h:4557 (averaging==2 || variance_form==residual selects the residual branch)

\subsection*{The members}

Table~\ref{tab:appendix-members} sets the six lattice members against the three switches, and it is
the equation-level counterpart of Table~\ref{tab:members} in the body: where that table says in
words what each member predicts, this one says which setting of Eqs.~\ref{eq:av0},
\ref{eq:variance-form} and~\ref{eq:mu-post} produces it.
% ADDED 2026-08-12 (09_carve_plan.md, reference check 3). This table was cited from nowhere at all,
% which is what let the body's Table 1 carry equation pointers into an appendix twenty pages later:
% two tables were doing one job and neither said so. The body table now carries English columns and
% names this one; this one names the equations.

% Non-floating: a float inside appendixbox's mdframed can escape the frame.
\begin{center}
\begin{tabularx}{\linewidth}{@{}lccl>{\raggedright\arraybackslash}X@{}}
\toprule
Member & window & recursive & interval variance & what it is \\
\midrule
\texttt{NR}  & an instant & no  & --- & instantaneous samples, windows independent \\
\texttt{R}   & an instant & yes & --- & instantaneous samples, posterior carried forward \\
\texttt{INR} & the start & no  & total & \texttt{NR} with the window average restored \\
\texttt{MR}  & the start & yes & total & \texttt{R} with the window average restored \\
\texttt{VR}  & the start & yes & residual & \texttt{MR} with switch~3 flipped, gain unchanged \\
\texttt{IR}  & both ends & yes & residual (forced) & \texttt{VR} with the boundary terms restored \\
\bottomrule
\end{tabularx}
\captionof{table}{The six members of the lattice as settings of the three switches. Reading down the
last column, each row differs from an earlier one in exactly one switch, which is what makes
the ladder measurable one step at a time. \texttt{IR} is MacroIR; \texttt{INR} is the
MacroINR of \citet{moffatt2025bayesian}. The data keys under which each is
dispatched are in Table~\ref{tab:family}.}
\label{tab:appendix-members}
\end{center}

One consequence of the two switches acting together is worth stating, because it is easy to
read the table as saying otherwise. Evaluated at the \emph{same} prior, \texttt{MR} and
\texttt{IR} assign the interval average the same total variance. Expanding the total form of
Eq.~\ref{eq:variance-form} and using
$\sum_{j}P_{i\to j}(\Delta)=1$ and $\sum_{j}P_{i\to j}(\Delta)\bar{\gamma}_{i\to j}
=\bar{\gamma}_{i\to\cdot}$,
\begin{equation}
  \bm{\mu}_0^\top\mathbf{w}\big|_{\text{total}}
  \;=\;
  \bm{\mu}_0^\top\mathbf{w}\big|_{\text{residual}}
  \;+\; \bm{\mu}_0^\top\mathbf{H}\mathbf{1}
  \;-\; \bar{\bm{\gamma}}_0^\top\mathrm{diag}(\bm{\mu}_0)\,\bar{\bm{\gamma}}_0,
  \label{eq:mr-ir-identity}
\end{equation}
and the two extra pieces are exactly what switch~1 moves from the third term of
Eq.~\ref{eq:pred-var} into the second when the window goes from the start state to both
ends. The
decompositions differ, the sum does not. What does change between the two members is the
gain, and at a common prior that is the whole of the \texttt{MR}-to-\texttt{IR} difference.
It does not follow that the two report the same variance on a recording. The predicted
variance divides the gain, so a member's variance sets the state the next window opens with:
the moment the gains differ the priors differ, and every quantity downstream of the prior,
the predicted variance included, differs with them. The identity above is between the two
maps from a prior to a predicted variance, and it is tested only where the two members still
hold the same prior.
% ADDED 2026-08-06. The paragraph asserted the identity without saying at what it is evaluated,
% which reads as a claim about a recording. On the figure-1 dumps MR and IR agree to the last
% printed digit at every interval where they still share a prior and separate at the first update
% after that. Source: papers/1_method/figures_build_plan.md:199-220, re-verified against
% legacy/qmodel.h:4568-4619 on 2026-08-06 (the boundary cross term enters gSg with + and ms with -
% and cancels; MR's contraction against the full P_Cov supplies exactly the piece the total form
% subtracts). Numbers are realization-dependent (the illustrative run is seeded with the random
% sentinel) and stay out of the body.

Two of the six pairs isolate one algebraic change at the point where it is made, and they are
the ones the Results use. Stepping \texttt{MR} to \texttt{VR} flips switch~3, and stepping
\texttt{VR} to \texttt{IR} restores the boundary terms in both the variance and the gain
while leaving the variance form where it already was. Neither step is confined to what it
flips, and the reason is the same one as above: the predictive variance divides the gain, so
flipping switch~3 moves the update as well, and the two steps are separable in the algebra
without being separable in what a recording does with them. The two together are the
\texttt{R}-to-\texttt{IR} gap.
% CORRECTED 2026-08-06. This paragraph carried the slogan "MR to VR flips switch 3 and touches
% nothing else", which papers/1_method/figures_build_plan.md:212, papers/_program/program.md:78,
% papers/1_method/decisions.md:161 and CONTINUE_HERE.md:88 all retire, in those documents' own
% words, because the predictive variance divides the gain: VR's smaller variance changes its update
% immediately and its propagated mean one step later. 05_discussion.tex:104 already states the
% qualification; this appendix contradicted it.

\subsection*{Least squares, which is off the lattice}

Classical nonlinear least squares has no setting of the three switches. It propagates the
mean occupancy alone, $\bm{\mu}_\Delta = \mathbf{P}(\Delta)^\top\bm{\mu}_0$, with no
covariance and no update, predicts $\bar{y}^{\mathrm{pred}}_{0\to\Delta}$ from Eq.~\ref{eq:ypred}, and
assigns one constant variance to every residual, fitted as the residual sum of squares over
the number of samples. Neither term of Eq.~\ref{eq:pred-var} that depends on the gating
survives, and there is no gain, so no occupancy object is ever conditioned on the data. It
is the anchor of the comparison rather than a member of the family.

\end{appendixbox}
